\chapter{定义一个值得训练的助手}\label{定义一个值得训练的助手}

\section{任务定义与评估 从一次失败开始}\label{任务定义与评估 从一次失败开始}

我们把一道题交给小助：``请解方程 3(a＋2)＝21，并解释每一步。''我们期待它先除以 3，再减去 2，得到 a＝5。它却接着写：``这是初中代数教材中的一道基础题，下面还有几道同类练习。''句子很通顺，也与题目有关，但它没有完成任务。

\emph{我们容易把这次失败归结为``模型不够聪明''。可如果把提示改成一段问答示例，再把同一道题放在末尾，它也许就能给出正确答案。}此时，算力和参数都没变，变化的是上下文。这提醒我们，至少要区分三个问题：模型能否表示所需知识；在当前提示下会生成什么；模型是否把这种生成转成了用户需要的行为。

基座语言模型学习的是文本延续规律。真实语料里，问题后面可能跟答案，也可能跟题目分析、引用、另一道题，甚至空白。我们给了一个任务，小助并不会自动地将这个任务作为唯一目标，因为它还不知道问题和回答之间的区别，哪里是问题的开始，哪里是问题的结束，哪里是回答的开始等。

在训练之前，我们通常先写一份小型验收约定。我们选出四类需求：讲清基础题；回答课程资料问题；在缺少资料时说不确定；通过计算器处理容易算错的数值。我们把``正确率''和``讲解质量''分开，把``引用存在''和``引用支持结论''分开。一个回答可以结果正确但过程错误，也可以格式漂亮却没有逻辑依据。

和传统机器学习一样，我们通常准备三份数据。训练集用来更新参数；验证集用来选提示、超参和检查点；最终测试集留到方案确定后使用。以后若反复根据同一份验证集修改方案，它事实上已经参与了开发，不能继续被当成完全独立的泛化证明。同一道母题的改写也应成组划分，避免训练和评估只是换了姓名与数字。当然，这是传统机器学习的思路，许多大模型的工程实践中基本没有验证集，研发人员直接合成训练数据并且在测试集上测试效果并且根据结果反馈调整训练集，尽管这是不够严谨的。

从这一步起，我们把``效果好像变好了''细化为成可核验的问题：哪类任务改善了，哪类退步了，在相同预算下是否仍成立？这份最初的评估约定，会在监督微调（Supervised Fine-Tuning，SFT）续训和强化学习（Reinforcement Learning，RL）过程中遇到麻烦时成为我们的定位工具。SFT 用示范回答训练模型，RL 则依据反馈优化模型的尝试；我们稍后会逐步展开。

\textbf{探索问题}~如何解释对于预训练模型：``\emph{我们容易把这次失败归结为``模型不够聪明''。可如果把提示改成一段问答示例，再把同一道题放在末尾，它也许就能给出正确答案。''}这个现象。

\textbf{本节相关资料}
\begin{enumerate}
    \item Brown 等，2020，\href{https://arxiv.org/abs/2005.14165}{Language Models are Few-Shot Learners} 。阅读重点：同一个模型为什么能通过上下文示例切换任务。
\end{enumerate}

\section{训练阶段 语言学习类比能解释什么}\label{训练阶段 语言学习类比能解释什么}

为了理解接下来发生的事，我们先想象一个熟悉的过程。婴儿长期听到周围的人说话，逐渐熟悉声音与表达；后来在互动中练习表达，得到示范和反馈。这能帮助我们建立一个入口：大量接触语言、学习具体表达、再依据结果调整，是彼此衔接但不相同的学习环节。

我们可以把预训练类比为广泛接触语言，把监督微调类比为有人示范``面对这个问题，可以这样回答''，把反馈学习类比为比较不同回答是否合适。这个类比让抽象的训练阶段有了生活背景，但仍然有一点区别：模型不会经历人类婴儿的发育过程。

婴儿有身体、感知、共同注意和社会互动。文本语言模型的预训练通常是在已有序列上预测下一个词元（token，即分词后的基本处理单元），并通过误差更新参数。注意，当我们提到大模型的行为时，经常指代的就是模型的``分布''。它在预训练时已经反复``尝试预测''，并非只听不说；SFT 也常使用同一种逐 token 交叉熵（cross-entropy）损失，并不是第一次获得生成能力。两者的重要区别在于数据的组织、监督的位置和我们希望强化的行为分布\footnote{理解大模型的输出（行为）是基于概率分布的，即根据上下文，从下一个token的概率分布中采样}。

预训练已经初步确定了 next token 分布，所以后训练叫做微调。不过因为注意力机制的存在，实际分布由参数和上下文共同决定，所以对于一个问题，后续解答时next token 分布不仅由问题本身决定，也跟其他上下文有关，比如问题前面可能还有一个解题指导，有指导地解题和无指导地解题难度肯定是不一样的，所以分布也有差异。全量后训练通常会训练大量数据，这样也可以比较明显改变模型的分布，这种规模的训练就不再是``微调''了。我们还可以做小规模后训练，比如在别人已经训练好的模型上继续用比较少的数据做SFT或者RL训练，这种小规模后训练的能力范围受训练起点、数据、优化与预算限制，通常很难大幅度调整模型的行为，甚至可能导致模型的通用能力收到影响。类似于一个人在大学系统性学习了四年中国哲学，后来自学了几个月西方哲学，那么西哲主要是做了一点思想补充和调整，其主干思维可能还是以中哲为主。

不妨想象小助对一道题能给出三种回答。训练前，它有时直接给答案，有时续写题目，有时解释步骤。SFT 提高给出高质量解题流程的概率；RL 依据奖励重新分配不同尝试的概率，比如同一个题目可能有好几种解题方案，可以通过给更好的解题方案更高的奖励来使得模型更倾向于生成这种更好的解题方案，这是一种偏好调整。若它从未形成某个关键运算过程，单靠稀疏奖励去生成完整解法可能很困难；若它已经偶尔做对，强化学习的奖励就可能帮助它更稳定地选到有效路径，比如解题 10 次，只对了2次，我们奖励这两次形成的解题轨迹，模型就会在这类问题上更倾向于走这条轨迹。

因此，这本书中的``学会''和``知道''都是方便交流的简写。落实到实验，我们总会追问：在哪些输入、怎样的解码条件和多少预算下，模型表现出了什么行为？有了这层约定，我们就能借助故事建立直觉，同时用机制纠正直觉。

\textbf{实践}~查资料或者实践看一下预训练数据和后训练数据有什么区别。

\textbf{本节相关资料}
\begin{enumerate}
    \item Ouyang 等，2022，\href{https://arxiv.org/abs/2203.02155}{Training language models to follow instructions with human feedback}，也就是 InstructGPT。阅读重点：如何利用人类反馈来优化大模型的行为。\textbf{该文章是人类偏好强化学习（RLHF）的里程碑论文之一，可以提前阅读，也可以在学到强化学习时阅读。}
\end{enumerate}
